\documentclass[student,print,12pt,a4paper]{physicsquiz}

\quiztitle{Starter Quiz Bank}
\quizcourse{Copyable local template}
\quizauthor{Course Team}
\quizversion{Practice}

\begin{document}
\makequiztitle

\begin{quizbank}
  \begin{quizquestion}[
    id=starter-units-001,
    topic=units,
    difficulty=foundation,
    marks=1,
    correct=B,
    tags={units,measurement},
    outcome={Identify a derived SI unit}
  ]
    Which unit is equivalent to a newton?
    \choices
      {\(\mathrm{kg\,m^{-1}\,s^{-2}}\)}
      {\(\mathrm{kg\,m\,s^{-2}}\)}
      {\(\mathrm{kg\,m^2\,s^{-1}}\)}
      {\(\mathrm{kg\,s^{-2}}\)}
      {\(\mathrm{m\,s^{-2}}\)}
  \end{quizquestion}
  \begin{quizsolution}
    A newton is the force needed to accelerate one kilogram by one metre per
    second squared, so the correct unit is \(\mathrm{kg\,m\,s^{-2}}\).
  \end{quizsolution}

  \begin{quizquestion}[
    id=starter-waves-001,
    topic=waves,
    difficulty=applied,
    marks=0.5,
    correct=C,
    tags={waves,frequency}
  ]
    A wave completes 20 oscillations in 5 seconds. What is its frequency?
    \choices
      {\(\SI{0.25}{Hz}\)}
      {\(\SI{2}{Hz}\)}
      {\(\SI{4}{Hz}\)}
      {\(\SI{15}{Hz}\)}
      {\(\SI{25}{Hz}\)}
  \end{quizquestion}
  \begin{quizsolution}
    Frequency is oscillations divided by time: \(20/5=\SI{4}{Hz}\).
  \end{quizsolution}
\end{quizbank}

\quizselectall

\begin{quizquestioncontent}
  \printquizquestions
\end{quizquestioncontent}

\begin{quizanswerkeycontent}
  \printquizanswerkey
\end{quizanswerkeycontent}

\begin{quizsolutioncontent}
  \printquizsolutions
\end{quizsolutioncontent}
\end{document}
